\documentclass[t,xcolor={usenames,dvipsnames}]{beamer}

\mode<presentation>
{
\usetheme{Warsaw}
%\setbeamercovered{transparent}
%\setbeamercolor{background canvas}{bg=white}
}

% Delete these, if you do not want the table of contents to pop up at
% the beginning of each (sub)section:
%\AtBeginSubsection[]
%{
%  \begin{frame}<beamer>{Outline}
%    \tableofcontents[currentsection,currentsubsection]
%  \end{frame}
%}
\AtBeginSection[]
{
  \begin{frame}<beamer>{Outline}
    \tableofcontents[currentsection]
  \end{frame}
}

\usepackage[english]{babel}
\usepackage[latin1]{inputenc}
\usepackage{times}
\usepackage[T1]{fontenc}
\usepackage{verbatim}
\usepackage{url}


% Author-date citations
\usepackage[authoryear,round]{natbib}
\let\cite=\citep  % default \cite such as {\LaTeX} authors are used to

% Where \includegraphics should look for figures
\graphicspath{{images/}}
\usepackage{epstopdf}
\DeclareGraphicsExtensions{.eps,.png,.jpg,.pdf}

\newcommand{\myhref}[2]{\href{#1}{\textcolor{Blue}{#2}}}
\newcommand{\subitem}[1]{\begin{itemize}\item #1 \end{itemize}}
\newcommand{\ghead}[1]{{\tiny #1\\}}

%%%%%%%%%%%%%%%%%%%%%%%%%%%%%%%%%%%%%%%%%%%%%%%%%%%%%%%%%%%%%%%%%%%%%%
% Title, author, logo
\title[Functional curation]{Functional curation of cardiac cell models}
\author{Jonathan Cooper and Gary Mirams}
\institute[University of Oxford]
{Computational Biology Group\\
 Oxford University Computing Laboratory}
\date{1\textsuperscript{st} February 2011}

\begin{document}

\begin{frame}
\titlepage
\end{frame}
%%%%%%%%%%%%%%%%%%%%%%%%%%%%%%%%%%%%%%%%%%%%%%%%%%%%%%%%%%%%%%%%%%%%%%

%\begin{frame}{Outline}
%\tableofcontents
%\end{frame}

\begin{frame}{What problem are we trying to solve?}
\begin{itemize}
\item Many (cardiac) cell models exist
\item Their properties are often not easy to determine
  \begin{itemize}
  \item May have to read several papers, and attempt to simulate the model
  \item This may not even be possible (at least without the author's assistance)
  \end{itemize}
% Waffle with some examples?
\item How can we determine a model's limitations or suitability for a given study?
\item When extending a model to explore a new experimental scenario, does it still reproduce the original behaviour?
\end{itemize}
\end{frame}

\begin{frame}{Example: S1-S2 protocol on canine models}
\begin{columns}[T]
\begin{column}{.33\linewidth}
\begin{center}
\ghead{\citealt{sicouri1991subpopulation}}\includegraphics[width=\textwidth]{sicouri_dog_ventricle_s1s2_curve}\\ \vspace{.1cm}
\only<2>{\ghead{Decker 2009}\includegraphics[width=\textwidth]{decker_2009_s1s2_curve}}
\end{center}
\end{column}
\begin{column}{.33\linewidth}
\begin{center}
\only<2>{\ghead{Fox 2002}\includegraphics[width=\textwidth]{fox_mcharg_gilmour_2002_s1s2_curve}\\ \vspace{.1cm}
\ghead{Hund--Rudy 2004}\includegraphics[width=\textwidth]{hund_rudy_2004_s1s2_curve}}
\end{center}
\end{column}
\begin{column}{.33\linewidth}
\begin{center}
\only<2>{\ghead{Lishitz--Rudy 2007}\includegraphics[width=\textwidth]{livshitz_rudy_2007_s1s2_curve}\\ \vspace{.1cm}
\ghead{Winslow 1999}\includegraphics[width=\textwidth]{winslow_model_1999_s1s2_curve}}
\end{center}
\end{column}
\end{columns}\end{frame}

\begin{frame}{What might we look for in a solution?}
\begin{itemize}[<+->]
\item A \alert{reference description} of the model in a standard format
  \subitem{CellML model repository}
\item A standard \alert{protocol language} to describe the simulations (and post-processing) performed
  \subitem{Need to create our own}
\item \alert{Automation} --- run all the simulations and display the results
  \subitem{Reference implementation based on Chaste}
\item Comparison with \alert{experimental data}
  \subitem{Initially `by eye'}
\end{itemize}
\end{frame}

\begin{frame}{Technical challenges}
\begin{itemize}
\item Dealing with variations in modelling conventions
  \begin{itemize}
  \item Need to run the same protocol on any model
    \subitem{e.g.\ variations in naming $\rightarrow$ annotation with \alert{ontologies}}
  \item Need to be able to compare models coherently
    \subitem{e.g.\ units conversions}
  \end{itemize}
\item Expressivity vs simplicity for the protocol language
  \begin{itemize}
  \item both for inputs
  \item and post-processing
  \end{itemize}
\end{itemize}
\end{frame}

\begin{frame}{Analysis challenges}
\begin{itemize}
\item Determining interesting protocols to run
\item Figuring out why models differ % Do we need an example based on the paper graphs?
\item Drawing insights into physiology
\end{itemize}
\end{frame}

\begin{frame}{Possible sub-projects}
\begin{itemize}
\item Computer Science
  \begin{itemize}
  \item Further developing the protocol language and implementation
  \item Optimising the system's performance
  \item Developing an online service to store, display, and compare results
  \end{itemize}
\item Modelling/Physiology
  \begin{itemize}
  \item Explore various different protocols, e.g. building on \citet{Romero*.09:Impact}
  \item Developing metrics to evaluate and display how well models replicate experimental data
  \end{itemize}
\item Both
  \begin{itemize}
  \item Developing and using ontologies to annotate the models
  \item Support for SBML to analyse reaction network models
  \end{itemize}
\end{itemize}
\end{frame}



%%%%%%%%%%%%%%%%%%%%%%%%%%%%%%%%%%%%%%%%%%%%%%%%%%%%%%%%%%%%%%%%%%%%%%
% Acknowledgments
%%%%%%%%%%%%%%%%%%%%%%%%%%%%%%%%%%%%%%%%%%%%%%%%%%%%%%%%%%%%%%%%%%%%%%

\begin{frame}{Acknowledgments}
\begin{columns}[T]
\begin{column}{.5\linewidth}
\begin{center}
\includegraphics[width=.3\textwidth]{LSI_hexagon_150mm}
\includegraphics[width=.3\textwidth]{SYSBIO_hexagon_150mm}\\ \vspace{.2cm}
\includegraphics[width=.9\textwidth]{chaste-266x60}
\end{center}
\end{column}
\begin{column}{.5\linewidth}
\begin{center}
\vspace{-1cm}
\includegraphics[width=.8\textwidth]{vph_noe_logo}\\
\includegraphics[width=.8\textwidth]{preDiCT}\\ \vspace{.2cm}
\includegraphics[width=.4\textwidth]{FP7-gen-RGB}
\end{center}
\end{column}
\end{columns}
\end{frame}

%%%%%%%%%%%%%%%%%%%%%%%%%%%%%%%%%%%%%%%%%%%%%%%%%%
\appendix
\begin{frame}[allowframebreaks]
\bibliographystyle{plainnat}
\bibliography{biblio}
\end{frame}

\end{document}
